\chapter{Einleitung}

\begin{itemize}
	\item welche inhaltlichen Schwerpunkte dabei gesetzt werden und welche Schritte zum Projektergebnis führen (sollen)
	\item kurze Begründung, warum dieses Thema gewählt wurde
	\item falls eine klare Arbeitsteilung in der Gruppe geplant ist, deren Aufteilung
	\item falls ein Anwendungsthema\footnote{z.\,B. aus einer anderen MINT-Wissenschaft oder mit einer externen Forschungseinrichtung, Firma etc.} bearbeitet wird, die Kooperation hier kurz erklären und die externe Betreuungsperson oder weitere betreuende Lehrperson des Fachgebietes angeben.
\end{itemize}


	Ziel unserer Projektarbeit war die systemorientierte Programmierung eines Computerspiels. \\
    Die Idee für die Art des Spiels kam aus einer gemeinsamen Erarbeitung. Zunächst entstanden zehn, ganz unterschiedliche, Spielideen, in der jeder der zehn Gruppenmitglieder einen Teil beigetragen hat. Anschließend wurde sich für eine Spielidee (später das \emph{CZGame}) entschieden. Diese Idee wurde nun weiter ausgereift. \\ 
	Aufgrund der Komplexität des Themas, wurde die Arbeit aufgeteilt und mehrere Gruppen entstanden. Diese Gruppen beschäftigten sich nun mit den unterschiedlichsten Themengebieten. Dabei gab es zu Beginn der Arbeit folgende Gruppen:
	\begin{enumerate}
		
		\item \textbf{Grafikdesign}
        
            Die Gruppe \emph{Grafikdesign}, bestehend aus Antonia, Jonas G. und Gustav, beschäftigte sich damit, die Hintergründe des Spiels zu gestalten. Da das Spiel im Pixelartdesign entstehen sollte , wurde für die Gestaltung der Hintergrundbilder das kostenlose Online-Tool \emph{Pixilart}  genutzt. Die Gruppe entwarf am Ende insgesamt 24 unterschiedliche Hintergründe, bestehend aus Räumen, Gängen und den Treppenhäusern.
	
		\item \textbf{Kampfkonzept}

            In der Gruppe \emph{Kampfkonzept} beschäftigten sich Sophie und Emil mit den einzelnen Abläufen der Handlungen, der im Spiel vorkommenden Charaktere. Besonders berücksichtigt wurde dabei die Interaktion zwischen der Lehrperson und der Spielerfigur. Am Ende entstand ein Konzept, bei dem Schüler und Lehrer abwechselnd versuchen sich ihrem gegenüber zu beweisen.
	
		\item \textbf{Minigames}
        
            Jonas L. und Rita beschäftigten sich in der Gruppe \emph{Minigames} mit dem Erstellen von kleinen Herausforderungen. Sie überlegten sich verschiedenste Konzepte für die verschieden Fächer, um dem Spiel eine tolle Abwechslung zu verpassen. Zudem erstellte die Gruppe jeweils drei verschiedene Schwierigkeitsstufen für jedes Minigame.  
	
		\item \textbf{Item- \& Charakterdesign}

            Wu-Yi und Arthur beschäftigten sich in der Gruppe \emph{Item- \& Charakterdesign} mit dem erstellen der nötigen Lehrpersonen, der Schüler und der Items. Hierfür sollte ebenfalls das Onlinetool \emph{Pixilart} genutzt werden.
	
		\item \textbf{Sound}

            Ashley erstellte Sounds und Musik, die während des Spielens im Hintergrund laufen sollten. Sie nutzte dafür ...
	
	\end{enumerate}

    Alle Gruppen beschäftigten sich zudem mit dem implementieren und programmieren ihrer Arbeitsergebnisse. Dabei erhielt jede Gruppe eine sehr große Unterstützung von Ashley, die ihr sehr großes Vorwissen bei der Programmierung einbrachte und sehr viel Zeit dafür investierte.
	
	 